\section{Background}
\label{sec:background}

\subsection{Speculative Decoding}
\label{sec:spec_decoding}

Autoregressive language models generate one token per forward pass, making inference latency proportional to output length.
Speculative decoding~\citep{ge2022lossless,pmlr-v202-leviathan23a,chen2023accelerating} accelerates the inference of a target model $M_t$ using a lightweight draft model $M_d$.
At each decoding cycle, the draft model proposes $\gamma$ candidate tokens $x_1, \ldots, x_{\gamma}$. The target model verifies all candidates in a single forward pass, accepting the longest prefix consistent with its own distribution.


Concretely, at each draft position $k$, the target model computes its own distribution $p_k^t$ and compares it against the draft distribution $p_k^d$.
The token $x_k$ is accepted with probability $\min(1,\, p_k^t(x_k) / p_k^d(x_k))$.
Verification proceeds left to right: the first rejection at position $k$ discards all subsequent tokens $x_{k+1}, \ldots, x_\gamma$, regardless of their quality.

Let $\tau$ denote the number of accepted tokens per cycle, and let $T_{\text{draft}}$ and $T_{\text{verify}}$ be the wall-clock times of the drafting and verification passes, respectively.
The average latency per generated token is:
\begin{equation}
L = \frac{T_{\text{draft}} + T_{\text{verify}}}{\tau}.
\label{eq:latency}
\end{equation}
Improving speedup therefore reduces to three levers: lowering $T_{\text{draft}}$ (draft faster), raising $\tau$ (draft better), or reducing the effective $T_{\text{verify}}$ (verify smarter).

\subsection{Drafter Architectures}
\label{sec:drafter_arch}

The design of the draft model determines how $T_{\text{draft}}$ and $\tau$ trade off.
Existing approaches fall into two categories.

\paragraph{Autoregressive drafters.}
Autoregressive drafters generate draft tokens sequentially, conditioning each position on previously sampled tokens~\citep{liu2024deepseek,pmlr-v235-li24bt,li-etal-2024-eagle,li2026eagle,zhang2025learningharmonizedrepresentationsspeculative}.
This explicit dependency gives strong modeling capacity, but the drafting cost grows linearly with block size: $T_{\text{draft}} \propto \gamma$, which forces autoregressive drafters to use small $\gamma$ and shallow architectures to keep $T_{\text{draft}}$ low.
To compensate for the short block, tree-based verification~\citep{miao2024specinfer} expands candidates into a tree and verifies multiple paths via tree attention, but the large number of verification tokens reduces overall serving throughput.

\paragraph{Parallel drafters.}
Parallel drafters produce all $\gamma$ draft tokens in a single forward pass, making $T_{\text{draft}}$ nearly independent of the block size~\citep{pmlr-v235-cai24b,chen2026dflash,liu2026dartdiffusioninspiredspeculativedecoding,li2025diffuspecunlockingdiffusionlanguage,sandler2025specdiff2scalingdiffusiondrafter}.
This allows substantially larger blocks (e.g., $\gamma{=}16$) without proportionally increasing latency.

Among them, DFlash~\citep{chen2026dflash} is a state-of-the-art parallel drafter, which conditions its draft model on rich context features extracted from the target model~(KV injection).
During prefill, hidden states from a set of target layers $\{l_1, \ldots, l_m\}$ are concatenated and projected into the draft hidden space:
\begin{equation}
H_{\text{ctx}} = \mathrm{RMSNorm}\bigl(W_c\,[H^{(l_1)};\, \ldots;\, H^{(l_m)}]\bigr),
\label{eq:ctx_feature}
\end{equation}
where $W_c \in \RR^{d \times md}$ is a shared projection.
These context features are injected into every draft layer by concatenating them with the draft block representations along the sequence dimension of keys and values:
\begin{equation}
K_i = [W_i^K H_{\text{ctx}};\; W_i^K H_d], \quad V_i = [W_i^V H_{\text{ctx}};\; W_i^V H_d].
\label{eq:kv_injection}
\end{equation}
All positions within a block attend bidirectionally to each other and to the injected target context.

The draft model shares the target model's embedding layer and language modeling head (both frozen).
It takes as input the embedding of an anchor token\footnote{We use the terms \textit{anchor token} and \textit{bonus token} interchangeably in this paper to denote the final token generated by the target model in the previous decoding round.} followed by $\gamma$ mask token embeddings, and produces logits for all mask positions in a single forward pass.
Since drafting requires only a single forward pass regardless of block size, DFlash can afford deeper architectures and larger blocks than autoregressive drafters under the same latency budget.